\begin{table}[htbp]\centering\small
\caption{Full horizon path for the external monetary-policy exposure specification}\label{tab:supp_external_horizons}
\begin{tabular}{rrrrr}\toprule
Horizon & $\hat\beta_h$ & 90\% CI lower & 90\% CI upper & $p$-value\\\midrule
0 & 0.062 & -0.088 & 0.212 & 0.487\\
1 & -0.102 & -0.407 & 0.204 & 0.576\\
2 & 0.091 & -0.230 & 0.412 & 0.633\\
3 & 0.113 & -0.155 & 0.381 & 0.478\\
4 & 0.272 & 0.108 & 0.436 & 0.009\\
5 & 0.181 & -0.057 & 0.419 & 0.206\\
6 & 0.134 & -0.070 & 0.338 & 0.273\\
7 & 0.097 & -0.109 & 0.304 & 0.430\\
8 & 0.190 & -0.158 & 0.538 & 0.360\\
9 & 0.384 & 0.085 & 0.684 & 0.037\\
10 & 0.252 & -0.050 & 0.554 & 0.167\\
11 & 0.057 & -0.357 & 0.471 & 0.817\\
12 & 0.035 & -0.331 & 0.402 & 0.871\\
\bottomrule\end{tabular}
\begin{singlespace}\footnotesize\textit{Notes:} The coefficient is the differential cumulative AERI response to a one-standard-deviation U.S. monetary-policy tightening for a country one standard deviation higher in predetermined 2014 external-debt exposure. Intervals and $p$-values use the finite-cluster reference with 28 denominator degrees of freedom.\end{singlespace}
\end{table}